\cleardoublepage
\phantomsection
\addcontentsline{toc}{chapter}{人工智能使用与审校声明}
\markboth{人工智能使用与审校声明}{人工智能使用与审校声明}
\FAFrontHeading{人工智能使用与审校声明}{使用边界与核验责任}

{\FAKai\color{FARose}\noindent\textbf{重要说明}\par}
\smallskip
本书不是经人类专家审定的正式教材、出版物或课程讲义. 数学学习具有累积性，任一处定义、定理假设、证明步骤或记号若存在错误，都可能影响后续理解；请将本书视为需要独立核验的人工智能生成型自学材料，而不是权威答案来源.

\FAFrontSubheading{制作方式}
本 PDF 的结构设计、正文生成、数学证明草拟、LaTeX 源码、跨章审计与排版均通过网页版 ChatGPT 的 Chat 功能完成；所用模型标识为 \textbf{GPT-5.6 Sol}，采用 \textbf{Extra High} 思考设置. 除项目发起者提出目标、约束、反馈与审美要求外，正文没有由人类作者逐条重写，证明也没有由人类专家逐项重新完成.

\FAFrontSubheading{人工检验状态}
本书\textbf{未经过独立人类数学专家的系统性逐条审校、同行评议或出版级事实核验}. 项目内部执行了依赖审计、证明闭包检查、记号与 LaTeX 规则检查、编译与视觉检查；这些步骤只能降低明显错误与排版问题的概率，\textbf{不能替代人工学术审查，也不构成正确性保证}.

\FAFrontSubheading{建议的核验方式}
若本书用于正式课程、考试、研究、论文写作或教学，至少应执行以下核验：
\begin{enumerate}[leftmargin=2.2em,itemsep=.35em,topsep=.35em]
  \item 对关键定理逐项核对数域、空间类别、拓扑、定义域与全部假设；
  \item 对负重证明独立复算关键估计、极限交换、紧性、可测性与定义域步骤；
  \item 与可靠教材或原始文献交叉核对谱论、无界算子、半群、Sobolev 空间、偏微分方程与非线性变分章节；
  \item 任何用于引用的内容都应回到可核验的正式来源，不应直接以本书作为唯一依据.
\end{enumerate}

\vfill
{\FAKai\small\color{FASlate}本声明用于透明说明人工智能生成方式与人工审校边界，使读者据此判断适用场景与所需核验强度.\par}
